%% =====================================================================
%%  Corrigés — Unit 19 « A trip to Boston »
%% =====================================================================
\corrigesUnite{19}{A trip to Boston}

\begin{corrige}{1}
\begin{reponses}[4]
\item \rep{went}
\item \rep{saw}
\item \rep{ate}
\item \rep{bought}
\item \rep{took}
\item \rep{had}
\item \rep{met}
\item \rep{thought}
\end{reponses}
\retiens{\en{bought} (\en{buy}) et \en{thought} (\en{think}) riment : \phon{bɑːt}, \phon{θɑːt}. Liste
complète page~\pageref{annexe-Verbes}.}
\end{corrige}

\begin{corrige}{2}
\begin{reponses}[3]
\item \rep{went}
\item \rep{took}
\item \rep{saw}
\item \rep{ate}
\item \rep{bought}
\item \rep{had}
\end{reponses}
\piege{\en{We goed}, \en{we taked}, \en{I buyed} n'existent pas : les verbes irréguliers ne prennent
jamais \en{-ed}.}
\end{corrige}

\begin{corrige}{3}
\begin{reponses}[1]
\item \rep{We didn't go to Harvard.}
\item \rep{Did she buy a cap?}
\item \rep{They didn't see the game.}
\item \rep{Did you take the train?}
\item \rep{He didn't eat lobster.}
\item \rep{Did Emma meet a student?}
\end{reponses}
\pourquoi{avec \en{did} et \en{didn't}, le verbe reprend sa base : \en{went} → \en{go}, \en{bought} →
\en{buy}, \en{saw} → \en{see}, \en{took} → \en{take}, \en{ate} → \en{eat}, \en{met} → \en{meet}.}
\end{corrige}

\begin{corrige}{4}
\begin{reponses}[1]
\item \rep{Where did you go?}
\item \rep{How long did you stay?}
\item \rep{How much did it cost?}
\item \rep{How did you travel?}
\item \rep{What did you see?}
\end{reponses}
\retiens{\en{How long…?} = durée ; \en{How much…?} = prix ; \en{How…?} = manière, moyen.}
\end{corrige}

\begin{corrige}{5}
\begin{reponses}[3][french]
\item \rep{V}
\item \rep{F}
\item \rep{V}
\item \rep{F}
\item \rep{V}
\item \rep{F}
\end{reponses}
\pourquoi{les colons refusaient la taxe ; ils ont jeté le thé à l'eau (\en{threw… into the water}) ; ce
fut le début de la révolution (\en{the beginning}).}
\end{corrige}

\begin{corrige}{6}
\begin{reponses}[6][french]
\item \rep{c}
\item \rep{a}
\item \rep{e}
\item \rep{b}
\item \rep{f}
\item \rep{d}
\end{reponses}
\end{corrige}

\begin{corrige}{7}
\begin{reponses}[3][french]
\item \rep{passé}
\item \rep{présent}
\item \rep{présent}
\item \rep{passé}
\item \rep{passé}
\item \rep{présent}
\end{reponses}
\pourquoi{\en{went}, \en{took}, \en{came} sont des prétérits ; \en{every day}, \en{at noon},
\en{on Sundays} accompagnent ici des habitudes au présent. Texte de la piste
page~\pageref{trans-19.10}.}
\end{corrige}

\begin{corrige}{8}
\begin{reponses}[5]
\item \rep{Saturday}
\item \rep{train}
\item \rep{four}
\item \rep{hotel}
\item \rep{lobster}
\end{reponses}
\pourquoi{texte de la piste page~\pageref{trans-19.11}.}
\end{corrige}

\begin{corrige}{9}
\begin{reponses}[4]
\item \rep{buy}
\item \rep{bought}
\item \rep{Did}
\item \rep{did}
\item \rep{took}
\item \rep{eat}
\item \rep{ate}
\item \rep{had}
\end{reponses}
\retiens{dans la question, \en{buy} et \en{eat} ; dans la réponse, \en{bought} et \en{ate}.}
\end{corrige}

\begin{corrige}{10}
\begin{reponses}[1]
\item \rep{We went to Boston last weekend.}
\item \rep{I took the train and I saw the sea.}
\item \rep{Did you buy souvenirs? — Yes, I did.}
\item \rep{She didn't eat lobster.}
\item \rep{How long did it take?}
\item \rep{Where did you stay?}
\item \rep{We met a nice student.}
\item \rep{It was amazing!}
\end{reponses}
\piege{phrase 4 : « elle n'a pas mangé » → \en{she didn't eat}, et non \en{she didn't ate}.}
\end{corrige}

\begin{corrige}{11}
Réponse modèle :
\begin{anglais}
\begin{tcolorbox}[otBase,colback=otGreenTint,colframe=otGreenTint,fontupper=\fontsize{10.4}{14}\selectfont]
Last summer, I went to Barcelona with my best friend. We took the plane from Paris. We stayed in a
small hotel near the beach. We saw the Sagrada Família and we walked in the old town. We ate tapas
every evening. It was amazing!
\end{tcolorbox}
\end{anglais}
\retiens{\en{take the plane} = « prendre l'avion » ; \en{my best friend} = « mon meilleur ami ».}
\end{corrige}

\begin{corrige}{12}
Réponses modèles :
\begin{reponses}[1]
\item \rep{I went to Rome.}
\item \rep{I took the plane.} \textit{ou} \rep{I took the train.}
\item \rep{I stayed in a hotel.} \textit{ou} \rep{I stayed with friends.}
\item \rep{I saw the Colosseum and a lot of churches.}
\item \rep{Yes, I did. I bought a T-shirt.} \textit{ou} \rep{No, I didn't.}
\end{reponses}
\end{corrige}

\begin{corrigetest}
\begin{reponses}[5][french]
\item \rep{c}
\item \rep{a}
\item \rep{b}
\item \rep{c}
\item \rep{a}
\item \rep{b}
\item \rep{a}
\item \rep{c}
\item \rep{b}
\item \rep{c}
\end{reponses}
\pourquoi{question 5 : \en{brought} est le prétérit de \en{bring} (apporter). Question 6 : \en{How much}
pour un prix, avec \en{did}.}
\end{corrigetest}
